\chapter{Analysis plan}
\label{chap:analysis}

The analysis is run after the session, on the committed captures, using the
application without modification. Running a modified pipeline on the
experimental data would mean the result does not describe the software that
exists.

\section{The pipeline, stage by stage}

The measurement path is the \emph{Geometric (Wire Tracking)} mode, which is
the application default. The alternative dense stereo mode is not used, and
Section \ref{sec:no-dense} explains why that matters more than it sounds.

\begin{enumerate}
    \item \textbf{Segmentation.} The cable detection tuner produces a binary
    mask for each image. Two segmenters are available: the classical one,
    based on a top-hat transform with adjustable parameters, and a trained
    random forest operating per pixel on nine hand-crafted features. Record
    which was used for each pair.

    \item \textbf{Endpoint detection.} The mask is skeletonised and pruned,
    and the two endpoints are taken as the geodesically farthest pair.

    \item \textbf{Path tracing.} Either automatic, a depth-first search over
    the mask with backtracking and momentum, or hand drawn by the operator
    clicking along the tether. Record which, per pair: this is what criterion
    A3.3 counts.

    \item \textbf{Rectification} of both images from the calibration.

    \item \textbf{Matching.} The point at $i\%$ along the left path is paired
    with the point at $i\%$ along the right path.

    \item \textbf{Triangulation} into a sparse point cloud, written out
    directly from the disparity and the baseline rather than through
    OpenCV's triangulation routines.

    \item \textbf{Metrics.} Arc length, straightness, depth range, diameter
    and curvature.

    \item \textbf{Export} to PLY, XYZ, PCD or OBJ.
\end{enumerate}

\section{What is computed, and against what}

\begin{table}[htbp]
\centering
\caption[Analysis outputs]{What is extracted from each analysed pair and what
it is compared against.}
\label{tab:analysis}
\begin{tabularx}{\textwidth}{@{}l X X@{}}
\toprule
\textbf{Quantity} & \textbf{Source} & \textbf{Compared against} \\
\midrule
Mask connectivity & stage 1 & criterion A3.1 \\
Endpoint count & stage 2 & criterion A3.2 \\
Tracking mode used, and coverage reached & stage 3 & criterion A3.3 \\
Retained sample fraction & stage 7 & criterion A3.4 \\
Arc length in metres & stage 7 & tape measure, criterion A3.5 \\
Straightness & stage 7 & criterion A3.6, on taut pairs only \\
Depth range & stage 7 & scene geometry record \\
Reported diameter & stage 7 & measured tether diameter \\
Curvature & stage 7 & qualitative, sag pairs \\
\bottomrule
\end{tabularx}
\end{table}

\section{The retained fraction has to be reported with every length}

\texttt{compute\_wire\_metrics} discards samples whose disparity is
non-positive or whose depth falls outside 0.05 to 50~m, and it does this
\emph{before} summing the chord lengths. A tether segment that gets filtered
out therefore shortens the reported arc length, and nothing in the output
says so. A length quoted without its retained fraction is not interpretable,
so the two are always reported together. Criterion A3.4 sets the floor at
90\%.

\section{The comparison that is expected to fail, and why that is useful}

Capture set S8 deliberately places one endpoint outside the frame of one
camera. The matching stage assumes both traced paths span the same physical
extent of the tether, and there is no epipolar check anywhere in the pipeline
that could detect a violation. The predicted consequence is a systematic
error in the reconstructed length with no accompanying warning.

The analysis of S8 is therefore: take the arc length reported for the S8
pairs, compare it against the tape-measured length and against the S1 pairs
of the same tether, and record the magnitude and the sign of the discrepancy.
This is the first time the project would have a measured size for an error it
has only ever reasoned about.

A secondary analysis, if the S8 result is large, is to run the same pairs
through the epipolar matcher that exists in the codebase but is not reachable
from the interface. That matcher does not rely on the parametric assumption
and was kept specifically against this possibility. It has never been
validated on real data, so a result from it is exploratory, not a
measurement.

\section{Why dense stereo is not used, and the trap that creates}
\label{sec:no-dense}

Dense stereo fails on a thin, textureless object because there is nothing
along the wire to correlate. This has already been demonstrated twice:
commercial structure-from-motion on the isolated tether returns a zero
resolution error and produces no camera alignment or sparse cloud at all, the
stated reason being insufficient visual information within the tether itself
rather than background interference; and dense stereo inside the laboratory
enclosure gives correct disparity on a solid target while also producing
non-physical disparity throughout regions that should be empty background,
yielding phantom geometry that needs manual intervention to clean
\cite{linku_d5_camera_report}.

The trap is operational. If the operator does not configure a mask in the
segmentation tuner, the processing falls back to dense stereo
\emph{automatically and silently}. It does not raise an error. It produces
something that looks like a point cloud. Any analysis run after skipping the
tuner is therefore running the method already known to fail, in a form that
looks plausible. The analysis log records, for every pair, that a mask was
present.

\section{Feeding the datasets}

Every pair that is successfully traced, automatically or by hand, is saved
into the path dataset together with its mask and its verified path. This
dataset currently holds three stereo samples against a recommended 30 to 100,
and it is the input to the one model in the plan that does not yet exist: a
learned per-branch scorer that would replace the fixed hand-tuned weights in
the tracker's decision at each crossing.

One caution applies to the other dataset. Saving a sample to the mask dataset
while the random-forest segmenter is active means training the model on its
own unverified output, which reinforces its errors. The application asks for
confirmation in that case; the answer during this session is no, unless the
mask has been inspected and corrected.
